\submittedby{Manas Singh}

Recall that the basis generated by a subbasis $\mathcal{S}$ is
$\B_{\mathcal{S}}=\{U_1\cap\dots\cap U_n\mid n\in\N,\ U_i\in\mathcal{S}\}$, and the topology
generated by $\mathcal{S}$ consists of all unions of elements of $\B_{\mathcal{S}}$.

\paragraph{(i)} Suppose $\mathcal{S}_1\subseteq\mathcal{S}_2$. Let $W\in\T_1$. Then
$W=\bigcup_{\alpha\in I}B_\alpha$, where each $B_\alpha$ is a finite intersection of elements of
$\mathcal{S}_1$. Since $\mathcal{S}_1\subseteq\mathcal{S}_2$, each set in that finite
intersection belongs to $\mathcal{S}_2$, so $B_\alpha\in\B_{\mathcal{S}_2}$ for all
$\alpha\in I$. Thus $W$ is a union of elements of $\B_{\mathcal{S}_2}$, i.e.\ $W\in\T_2$.
Therefore $\T_1\subseteq\T_2$.

\paragraph{(ii)} We give an explicit example where $\T_1=\T_2$ even though
$\mathcal{S}_1\neq\mathcal{S}_2$.

Let $X=\{1,2,3\}$ and define
\[ \mathcal{S}_1=\bigl\{\{1\},\{2\},X\bigr\},\qquad
   \mathcal{S}_2=\bigl\{\{1\},\{2\},\{1,2\},X\bigr\}. \]
Both cover $X$ (they contain $X$), so both are subbases. Clearly
$\mathcal{S}_1\neq\mathcal{S}_2$, since $\{1,2\}\in\mathcal{S}_2\setminus\mathcal{S}_1$.
\begin{itemize}
  \item The finite intersections of elements of $\mathcal{S}_1$ are $\{1\}$, $\{2\}$, $X$ and
  $\{1\}\cap\{2\}=\emptyset$, so $\B_1=\{\emptyset,\{1\},\{2\},X\}$. Taking all unions (note
  $\{1\}\cup\{2\}=\{1,2\}$),
  \[ \T_1=\{\emptyset,\{1\},\{2\},\{1,2\},X\}. \]
  \item The finite intersections of elements of $\mathcal{S}_2$ are $\{1\},\{2\},\{1,2\},X$ and
  $\emptyset$ (for instance $\{1\}\cap\{2\}$), and no new sets arise from intersecting these
  (e.g.\ $\{1,2\}\cap\{1\}=\{1\}$). So $\B_2=\{\emptyset,\{1\},\{2\},\{1,2\},X\}$, and since this
  family is already closed under unions,
  \[ \T_2=\{\emptyset,\{1\},\{2\},\{1,2\},X\}. \]
\end{itemize}
Thus $\T_1=\T_2$ although $\mathcal{S}_1\neq\mathcal{S}_2$.

\submittedby{Sarthak}

\paragraph{(i) Proof that $\T_1\subseteq\T_2$.} By definition, the topology generated by a
subbasis $\mathcal{S}$ is the smallest topology on $X$ containing $\mathcal{S}$. Equivalently,
it is the collection of arbitrary unions of finite intersections of elements of $\mathcal{S}$.

Since $\mathcal{S}_1\subseteq\mathcal{S}_2$, and since $\T_2$ is a topology containing
$\mathcal{S}_2$, we have
\[ \mathcal{S}_1\subseteq\mathcal{S}_2\subseteq\T_2. \]
Thus $\T_2$ is a topology on $X$ that contains $\mathcal{S}_1$. But $\T_1$ is the smallest
topology on $X$ containing $\mathcal{S}_1$. Therefore,
\[ \T_1\subseteq\T_2. \]

Alternatively, let $U\in\T_1$. Then $U$ can be written as a union of finite intersections of
elements of $\mathcal{S}_1$:
\[ U=\bigcup_{\alpha\in A}\Bigl(\bigcap_{i=1}^{n_\alpha}S_{\alpha i}\Bigr), \]
where each $S_{\alpha i}\in\mathcal{S}_1$. Since $\mathcal{S}_1\subseteq\mathcal{S}_2$, each
$S_{\alpha i}\in\mathcal{S}_2$. Hence each finite intersection
\[ \bigcap_{i=1}^{n_\alpha}S_{\alpha i} \]
belongs to $\T_2$. Since $\T_2$ is closed under arbitrary unions, it follows that
\[ U\in\T_2. \]
Therefore,
\[ \boxed{\T_1\subseteq\T_2.} \]

\paragraph{(ii) Example where $\T_1=\T_2$.} Let
\[ X=\R, \]
and define
\[ \mathcal{S}_1=\{(a,b) : a<b\}, \]
the collection of all open intervals in $\R$, and
\[ \mathcal{S}_2=\{(a,b) : a<b\}\cup\{(a,\infty) : a\in\R\}. \]
Clearly,
\[ \mathcal{S}_1\subsetneq\mathcal{S}_2, \]
because each ray $(a,\infty)$ is in $\mathcal{S}_2$ but is not an open interval of the form
$(c,d)$.

Let $\T_1$ be the topology generated by $\mathcal{S}_1$, and let $\T_2$ be the topology
generated by $\mathcal{S}_2$. The collection $\mathcal{S}_1$ generates the usual topology on
$\R$. Also, every ray is a union of open intervals:
\[ (a,\infty)=\bigcup_{n=1}^{\infty}(a,a+n). \]
Hence $(a,\infty)\in\T_1$, so
\[ \mathcal{S}_2\subseteq\T_1. \]
Therefore, since $\T_2$ is the smallest topology containing $\mathcal{S}_2$, we get
\[ \T_2\subseteq\T_1. \]
On the other hand, $\mathcal{S}_1\subseteq\mathcal{S}_2$, so by part (i),
\[ \T_1\subseteq\T_2. \]
Combining both inclusions,
\[ \T_1=\T_2. \]
Thus, even though $\mathcal{S}_1\subsetneq\mathcal{S}_2$, both subbases generate the same
topology, namely the usual topology on $\R$.
